\paragraph{Current-state timing selection and transactional repair.}
We now test an external automatic repair loop from the unmodified frozen Tf and strong native R1 states. At each accepted state, it ranks all eligible routed connections using the newly computed native setup slack, groups numerical ties within eight float ULPs, and queries full-device fixed-tree alternatives. This tolerance groups targets only; it is never an acceptance tolerance. Single-connection candidates retain every other branch and other-net occupancy. When those candidates do not advance the state, conflict-based joint path search reserves compatible target paths and restores displaced sinks sequentially under the updated occupancy. The complete transaction receives independent physical checks and native FULL/STA before its immutable state pointer can be committed. Rejected transactions retain the original state. State/candidate identities suppress repeated unchanged proposals; equivalent historical queries and complete transactions are explicitly reused.

\paragraph{Exact endpoint policy and branch divergence.}
The requested \texttt{eltwise\_layer.sdc} was absent in these runs. We therefore export the actual native default constraints, including \texttt{virtual\_io\_clock}, \texttt{clk}, and \texttt{clk\_mem}; paths between the two netlist-clock domains are ignored, while the applicable virtual-I/O domain pairs remain. We do not introduce new constraints retrospectively. The fixed timing-graph SINK set contains 13,744 endpoints and 14,032 setup endpoint/domain components. The strict policy requires a strictly smaller native least-slack critical-path delay. The alternative requires nonincreasing CPD and componentwise nondecreasing setup slack with at least one strict improvement, over the same fixed endpoint/domain keys and applied constraints. Both require legal routing and no increase in negative hold violations. Native values are exported at full precision and compared without epsilon, preventing accumulated tolerated regression. Positive hold margin can decrease while remaining nonviolating. Unconstrained identities and nonfinite status must remain unchanged.

The earlier legal single-connection R1 plateau and the new A0 coordinated plateau change physical connection delays but leave the entire endpoint setup vector unchanged; both policies reject them. From original R1, the automatic joint proposal for nets 6587/6590 pin 4 reduces CPD from 6.018101129 to 5.854659868 ns and improves 246 setup components without regression. Both branches accept it and independently reselect nets 76/48 pin 1; neither has a shorter fixed-background capacity-feasible connection. This automatically generated joint route uses 41 SB pairs and 971 CB source resources, versus 39 and 969 at the start. It matches the earlier manually developed CPD but uses a different blocker assignment; the earlier manual route's smaller resource counts must not be substituted.

From original Tf, the strict branch accepts current-selected repairs of net 5479/pin 36, coordinated net 6587/pin 4, net 7647/pin 56, and net 6583/pin 24. CPD progresses through 6.203926706, 6.096180893, 6.019283294, and 6.018101129 ns. The final state improves 256 setup components but regresses one relative to original Tf, and uses 130 SB pairs, 961 CB sources, and planar wirelength 228511, versus 118, 955, and 228420 initially. The endpoint branch accepts only the first step: the second improves 254 components but worsens \texttt{out:PRDATA[7]} in the \texttt{clk}-to-virtual-I/O domain by 0.460476435 ns. It therefore remains at 6.203926706 ns and independently rejects that transaction. There is still no demonstrated strict-reject/endpoint-accept case. At the strict branch's later state, a single-net plateau is rejected, followed by a legal joint candidate that worsens CPD to 6.203926706 ns and is rolled back. These observations motivate timing-aware treatment of displaced connections; a local target-delay optimum is not a circuit timing optimum.

\paragraph{Failure, scope and cost accounting.}
A separate initial four-target Tf joint proposal restores 78 displaced sinks but worsens CPD from 6.208121572 to 6.734843794 ns; it is preserved as a legal rejected candidate. Its recorded construction time is 549.936 s, including 91.053 s of query subprocess time, but excludes the earlier input-binding stage. The native evaluation wrapper takes 40.433 s, including 8.187 s of VPR execution. The missing earlier cost is explicitly unknown. The R1 and Tf two-policy orchestrators take 125.741 and 361.208 s, respectively, with shared, explicitly reused evidence; these are neither independent per-arm timings nor cold-start complete-method measurements. The new ten-sink Tf transaction reuses ten exact query inputs and takes 36.326 s for the complete wrapper, dominated by state construction and checking. Historical U formal routing times are not used for this prototype. The next integration decision should follow these measured preparation/audit costs. These remain one-circuit development results, not a robust method promotion or a stopping condition for the research.

\paragraph{Matched native connection-search control.}
An opt-in control calls VPR's actual connection router and \texttt{update\_from\_heap} on copied, pruned trees. Twelve initial native searches match the external operator's retained nodes, expandable anchors and full-device domain. Native transformed criticality with congestion pressure 1 returns capacity-conflicting improving paths for the two R1 targets. Continuing those same targets at pressure 10 produces legal paths identical to the capacity-constrained operator, including every net's edge, switch and terminal set; each improves its physical connection delay by 0.185825799 ns. Both remain full-precision endpoint and CPD plateaus and are rejected by both policies. A native equal-delay alternative for net 6639/pin 1 is also legal but rejected. This is evidence against attributing these single-connection results to a distinct query mechanism. The following matched joint-native controls examine recovery and scheduling before attributing the joint CPD gain. One initial native-control integration failure used a stale post-prune sink pointer for a diagnostic assertion; it is preserved and corrected using RR-tree membership. The completed control uses 14 native connection searches, despite the inherited wrapper field denoting zero whole \texttt{--route} invocations.
